\documentclass[../libro]{subfiles}

\begin{document}

\ifSubfilesClassLoaded{\mainmatter\chapter{行列式}\clearpage}{}

\KunAsteriskoEnEnhavtabelo
\section{Binet--Cauchy 公式}
\SenAsteriskoEnEnhavtabelo

\maldevigalegajxo

或许, (方) 阵的行列式与阵的积的定义是较复杂的,
跟阵的转置、加、减、数乘比.
不过, 它们有不平凡的关系.

设 \(A\) 是一个 \(m \times n\)~阵,
且 \(B\) 是一个 \(n \times m\)~阵.
根据阵的积的定义,
\(BA\) 是一个 \(n\)~级阵,
故它有行列式.

\begin{theorem}[Binet--Cauchy 公式]
    设 \(A\), \(B\) 分别是 \(m \times n\), \(n \times m\)~阵.

    (1)
    若 \(n > m\), 则 \(\det {(BA)} = 0\).

    (2)
    若 \(n \leq m\),
    则
    \begin{align*}
             & \det {(BA)}
        \\
        = {} & \sum_{1 \leq j_1 < j_2 < \dots < j_n \leq m}
        {
            \det {\left(
                B\binom{1, 2, \dots, n}{j_1, j_2, \dots, j_n}
                \right)}
            \det {\left(
                A\binom{j_1, j_2, \dots, j_n}{1, 2, \dots, n}
                \right)}
        }.
    \end{align*}
    特别地, 若 \(n = m\), 则因适合条件
    \(1 \leq j_1 < j_2 < \dots < j_n \leq m\)
    的 \(j_1\), \(j_2\), \(\dots\), \(j_n\),
    是, 且只能是,
    \(1\), \(2\), \(\dots\), \(n\),
    故
    \begin{align*}
        \det {(BA)}
        = {} &
        \det {\left(
            B\binom{1, 2, \dots, n}{1, 2, \dots, n}
            \right)}
        \det {\left(
            A\binom{1, 2, \dots, n}{1, 2, \dots, n}
            \right)}
        \\
        = {} & \det {(B)} \det {(A)}.
    \end{align*}
\end{theorem}

我用二个例助您理解, 此定理在说什么.

\begin{example}
    设
    \begin{align*}
        A =
        \begin{bmatrix}
            1 & 3 & 5 \\
            2 & 4 & 6 \\
        \end{bmatrix},
        \quad
        B =
        \begin{bmatrix}
            7 & 10 \\
            8 & 11 \\
            9 & 12 \\
        \end{bmatrix}.
    \end{align*}
    不难算出,
    \begin{align*}
        AB =
        \begin{bmatrix}
            76  & 103 \\
            100 & 136 \\
        \end{bmatrix},
        \quad
        BA =
        \begin{bmatrix}
            27 & 61 & 95  \\
            30 & 68 & 106 \\
            33 & 75 & 117 \\
        \end{bmatrix}.
    \end{align*}
    直接计算, 可知
    \begin{align*}
        \det {(AB)} = 76 \cdot 136 - 100 \cdot 103 = 36,
    \end{align*}
    而
    \begin{align*}
        \det {(BA)}
        = {} & \hphantom{{} + {}}
        27 \cdot 68 \cdot 117
        + 30 \cdot 75 \cdot 95
        + 33 \cdot 61 \cdot 106
        \\
             &
        - 27 \cdot 75 \cdot 106
        - 30 \cdot 61 \cdot 117
        - 33 \cdot 68 \cdot 95
        \\
        = {} & 0.
    \end{align*}

    或许, 前面的计算是较复杂的.
    我们试用 Binet--Cauchy 公式, 再计算它们.

    因为 \(A\), \(B\) 的尺寸分别是 \(2 \times 3\), \(3 \times 2\),
    而 \(3 > 2\),
    根据 (1), \(\det {(BA)} = 0\)
    (注意, \(BA\) 是 \(3\)~级阵).

    因为 \(2 \leq 3\), 故, 根据 (2),
    \begin{align*}
             & \det {(AB)}
        \\
        = {} &
        \hphantom{{} + {}}
        \det {\left(A\binom{1,2}{1,2}\right)}
        \det {\left(B\binom{1,2}{1,2}\right)}
        + \det {\left(A\binom{1,2}{1,3}\right)}
        \det {\left(B\binom{1,3}{1,2}\right)}
        \\
             &
        + \det {\left(A\binom{1,2}{2,3}\right)}
        \det {\left(B\binom{2,3}{1,2}\right)}.
    \end{align*}
    这里, 注意, 适合条件
    ``\(1 \leq j_1 < j_2 \leq 3\)''
    的 \((j_1, j_2)\)
    恰有三个:
    \((1, 2)\), \((1, 3)\), \((2, 3)\).
    不难写出
    \begin{align*}
         & A\binom{1,2}{1,2}
        = \begin{bmatrix} 1&3\\2&4\\ \end{bmatrix},
        \quad
        B\binom{1,2}{1,2}
        = \begin{bmatrix} 7&10\\8&11\\ \end{bmatrix}; \\
         & A\binom{1,2}{1,3}
        = \begin{bmatrix} 1&5\\2&6\\ \end{bmatrix},
        \quad
        B\binom{1,3}{1,2}
        = \begin{bmatrix} 7&10\\9&12\\ \end{bmatrix}; \\
         & A\binom{1,2}{2,3}
        = \begin{bmatrix} 3&5\\4&6\\ \end{bmatrix},
        \quad
        B\binom{2,3}{1,2}
        = \begin{bmatrix} 8&11\\9&12\\ \end{bmatrix}.
    \end{align*}
    从而
    \begin{align*}
        \det {(AB)}
        = (-2) \cdot (-3) + (-4) \cdot (-6) + (-2) \cdot (-3)
        = 36.
    \end{align*}
\end{example}

\begin{example}
    设
    \begin{align*}
        A = \begin{bmatrix}
                3  & 2 \\
                10 & 7 \\
            \end{bmatrix},
        \quad
        B = \begin{bmatrix}
                5  & 4 \\
                11 & 9 \\
            \end{bmatrix}.
    \end{align*}
    不难算出
    \begin{align*}
        AB
        = \begin{bmatrix}
              37  & 30  \\
              127 & 103 \\
          \end{bmatrix},
        \quad
        BA
        = \begin{bmatrix}
              55  & 38 \\
              123 & 85 \\
          \end{bmatrix}.
    \end{align*}
    可以看到, \(AB \neq BA\).

    不过, \(\det {(AB)} = \det {(BA)}\).
    一方面, 我们可直接验证:
    \begin{align*}
         & \det {(AB)} = 37 \cdot 103 - 127 \cdot 30 = 1, \\
         & \det {(BA)} = 55 \cdot 85 - 123 \cdot 38 = 1.
    \end{align*}
    另一方面, Binet--Cauchy 公式指出,
    \begin{align*}
        \det {(AB)}
        = {} & \det {(A)} \det {(B)} \\
        = {} & \det {(B)} \det {(A)} \\
        = {} & \det {(BA)}.
    \end{align*}
    毕竟, 数的乘法是可换的.
\end{example}

为简单地论证公式,
我们要三个新事实;
为证明这三个新事实的前二个,
我们要一些老的公式.

% subfile fix: repeat restatable
% see:
% determinanto - determinantoj.tex
\ifSubfilesClassLoaded{%
    \begin{restatable}{theorem}{TheoremFullExpansion}
        设 \(A\) 是 \(n\)~级阵 (\(n \geq 1\)).
        设 \(j_1\), \(j_2\), \(\dots\), \(j_n\) 是%
        不超过 \(n\) 的正整数,
        且互不相同.
        则
        \begin{align*}
                 & \det {(A)}
            \\
            = {} &
            \sum_{\substack{
            1 \leq i_1, i_2, \dots, i_n \leq n \\
                    i_1, i_2, \dots, i_n\,\text{互不相同}
                }}
            {s(i_1, i_2, \dots, i_n)\,
                s(j_1, j_2, \dots, j_n)\,
                [A]_{i_1,j_1} [A]_{i_2,j_2} \dots [A]_{i_n,j_n}}.
        \end{align*}
        特别地, 取 \(j_1\), \(j_2\), \(\dots\), \(j_n\)
        为 \(1\), \(2\), \(\dots\), \(n\),
        并注意, \(s(1, 2, \dots, n) = 1\),
        得
        \begin{align*}
            \det {(A)}
            = {} &
            \sum_{\substack{
            1 \leq i_1, i_2, \dots, i_n \leq n \\
                    i_1, i_2, \dots, i_n\,\text{互不相同}
                }}
            {s(i_1, i_2, \dots, i_n)\,
                [A]_{i_1,1} [A]_{i_2,2} \dots [A]_{i_n,n}}.
        \end{align*}
    \end{restatable}
}{\TheoremFullExpansion*}

% subfile fix: repeat restatable
% see:
% determinanto - determinantoj.tex
\ifSubfilesClassLoaded{%
    \begin{restatable}{theorem}{TheoremColumnSwapAndSign}
        设 \(n\)~级阵 \(A\) 的列~\(1\), \(2\), \(\dots\), \(n\)
        分别为 \(a_1\), \(a_2\), \(\dots\), \(a_n\).
        设 \(\ell_1\), \(\ell_2\), \(\dots\), \(\ell_n\) 是%
        不超过 \(n\) 的正整数,
        且互不相同.
        则
        \begin{align*}
            \det {[a_{\ell_1}, a_{\ell_2}, \dots, a_{\ell_n}]}
            = s(\ell_1, \ell_2, \dots, \ell_n) \det {(A)}.
        \end{align*}
    \end{restatable}
}{\TheoremColumnSwapAndSign*}

现在, 我要开始展现三个新事实了.

\begin{theorem}
    设 \(C\) 为 \(n \times m\)~阵,
    且 \(m \geq n\).
    设 \(C\)~的列~\(1\), \(2\), \(\dots\), \(m\)
    分别为 \(c_1\), \(c_2\), \(\dots\), \(c_m\).
    设 \(k_1\), \(k_2\), \(\dots\), \(k_n\) 是%
    不超过 \(m\)~的正整数,
    且互不相同.
    设 \(j_1\), \(j_2\), \(\dots\), \(j_n\) 是
    \(k_1\), \(k_2\), \(\dots\), \(k_n\) 的自然排列.
    则
    \begin{align*}
        \det {[c_{k_1}, c_{k_2}, \dots, c_{k_n}]}
        = s(k_1, k_2, \dots, k_n)
        \det {[c_{j_1}, c_{j_2}, \dots, c_{j_n}]}.
    \end{align*}
\end{theorem}

\begin{proof}
    设 \(k_t\) 在 \(k_1\), \(k_2\), \(\dots\), \(k_n\)
    的自然排列
    \(j_1\), \(j_2\), \(\dots\), \(j_n\)
    里的位置为 \(f(k_t)\)
    (\(t = 1\), \(2\), \(\dots\), \(n\)).
    注意, \(f(j_t) = t\).
    并且, 若 \(k_t < k_v\), 必有 \(f(k_t) < f(k_v)\).
    反过来, 若 \(f(k_t) < f(k_v)\), 必有 \(k_t < k_v\).
    (我用一个简单的例助您理解.
    设 \(m = 9\), \(n = 3\).
    设 \(k_1\), \(k_2\), \(k_3\) 分别为 \(8\), \(9\), \(6\).
    那么, \(k_1\), \(k_2\), \(k_3\) 的自然排列
    \(j_1\), \(j_2\), \(j_3\)
    是
    \(6\), \(8\), \(9\).
    故 \(f(k_1) = f(j_2) = f(8) = 2\),
    \(f(k_2) = f(j_3) = f(9) = 3\),
    且 \(f(k_3) = f(j_1) = f(6) = 1\).)
    从而
    \(\operatorname{sgn} {(k_v - k_t)}
    = \operatorname{sgn} {(f(k_v) - f(k_t))}\).

    作阵~\(G = [c_{j_1}, c_{j_2}, \dots, c_{j_n}]\).
    设 \(G\) 的列~\(1\), \(2\), \(\dots\), \(n\)
    分别是
    \(g_1\), \(g_2\), \(\dots\), \(g_n\).
    则
    \begin{align*}
        [c_{k_1}, c_{k_2}, \dots, c_{k_n}]
        = [g_{f(k_1)}, g_{f(k_2)}, \dots, g_{f(k_n)}].
    \end{align*}
    故
    \begin{align*}
        \det {[c_{k_1}, c_{k_2}, \dots, c_{k_n}]}
        = {} &
        \det {[g_{f(k_1)}, g_{f(k_2)}, \dots, g_{f(k_n)}]}
        \\
        = {} &
        s(f(k_1), f(k_2), \dots, f(k_n))
        \det {[g_1, g_2, \dots, g_n]}
        \\
        = {} &
        s(k_1, k_2, \dots, k_n)
        \det {[c_{j_1}, c_{j_2}, \dots, c_{j_n}]}.
        \qedhere
    \end{align*}
\end{proof}

\begin{theorem}
    设 \(D\) 为 \(m \times n\)~阵,
    且 \(m \geq n\).
    设 \(i_1\), \(i_2\), \(\dots\), \(i_n\) 是%
    不超过 \(m\) 的正整数,
    且 \(i_1 < i_2 < \dots < i_n\).
    则
    \begin{align*}
             &
        \det {\left(
            D\binom{i_1, i_2, \dots, i_n}{1, 2, \dots, n}
            \right)}
        \\
        = {} &
        \sum_{\substack{
        k_1, k_2, \dots, k_n\,\text{是} \\
                i_1, i_2, \dots, i_n\,\text{的排列}
            }}
        {s(k_1, k_2, \dots, k_n)\,
            [D]_{k_1,1} [D]_{k_2,2} \dots [D]_{k_n,n}}.
    \end{align*}
\end{theorem}

\begin{proof}
    设
    \begin{align*}
        H =
        D\binom{i_1, i_2, \dots, i_n}{1, 2, \dots, n}.
    \end{align*}
    不难看出,
    \([H]_{u,v} = [D]_{i_u,v}\).

    记 \(f(i_t) = t\)
    (\(t = 1\), \(2\), \(\dots\), \(n\)).
    设 \(k_1\), \(k_2\), \(\dots\), \(k_n\)
    是 \(i_1\), \(i_2\), \(\dots\), \(i_n\)
    的一个排列.
    那么, \(f(k_1)\), \(f(k_2)\), \(\dots\), \(f(k_n)\)
    是 \(1\), \(2\), \(\dots\), \(n\) 的一个排列.
    反过来, 若 \(v_1\), \(v_2\), \(\dots\), \(v_n\) 是
    \(1\), \(2\), \(\dots\), \(n\) 的一个排列,
    则 \(i_{v_1}\), \(i_{v_2}\), \(\dots\), \(i_{v_n}\)
    是 \(i_1\), \(i_2\), \(\dots\), \(i_n\) 的一个排列.
    并且, 若 \(k_t < k_v\), 必有 \(f(k_t) < f(k_v)\).
    反过来, 若 \(f(k_t) < f(k_v)\), 必有 \(k_t < k_v\).
    从而
    \(\operatorname{sgn} {(k_v - k_t)}
    = \operatorname{sgn} {(f(k_v) - f(k_t))}\).
    故
    \begin{align*}
             &
        \det {\left(
            D\binom{i_1, i_2, \dots, i_n}{1, 2, \dots, n}
            \right)}
        \\
        = {} &
        \det {(H)}
        \\
        = {} &
        \sum_{\substack{
        k_1, k_2, \dots, k_n\,\text{是} \\
                i_1, i_2, \dots, i_n\,\text{的排列}
            }}
        {s(f(k_1), f(k_2), \dots, f(k_n))\,
            [H]_{f(k_1),1} [H]_{f(k_2),2} \dots [H]_{f(k_n),n}}
        \\
        = {} &
        \sum_{\substack{
        k_1, k_2, \dots, k_n\,\text{是} \\
                i_1, i_2, \dots, i_n\,\text{的排列}
            }}
        {s(k_1, k_2, \dots, k_n)\,
            [D]_{k_1,1} [D]_{k_2,2} \dots [D]_{k_n,n}}.
        \qedhere
    \end{align*}
\end{proof}

\begin{theorem}
    设 \(n \leq m\).
    设数 \(f(k_1, k_2, \dots, k_n)\)
    (其中
    \(k_1\), \(k_2\), \(\dots\), \(k_n\) 过
    \(1\), \(2\), \(\dots\), \(m\))
    适合:
    若 \(k_p = k_q\) (\(1 \leq p < q \leq n\)),
    则 \(f(k_1, k_2, \dots, k_n) = 0\).
    那么
    \begin{align*}
             &
        \sum_{1 \leq k_1, k_2, \dots, k_n \leq m}
        {f(k_1, k_2, \dots, k_n)}
        \\
        = {} &
        \sum_{\substack{
        1 \leq k_1, k_2, \dots, k_n \leq m \\
                k_1, k_2, \dots, k_n\,\text{互不相同}
            }}
        {f(k_1, k_2, \dots, k_n)}
        \\
        = {} &
        \sum_{1 \leq j_1 < j_2 < \dots < j_n \leq m}
        {\sum_{\substack{
        k_1, k_2, \dots, k_n\,\text{是}     \\
                    j_1, j_2, \dots, j_n\,\text{的排列}
                }}
            {f(k_1, k_2, \dots, k_n)}}.
    \end{align*}
\end{theorem}

\begin{proof}
    由 \(f\) 的性质可知, 前一个等式成立.
    后一个等式也不难.
    注意,
    要从 \(m\)~个不同的数里有前后次序地选
    \(n\)~个不同的数
    \(k_1\), \(k_2\), \(\dots\), \(k_n\),
    我们可以%
    先从小到大地从这 \(m\)~个数里选 \(n\)~个%
    不同的数
    \(j_1\), \(j_2\), \(\dots\), \(j_n\),
    且再有前后次序地作这 \(n\)~个数的排列
    \(k_1\), \(k_2\), \(\dots\), \(k_n\).
    所以, 后一个等式也成立.
\end{proof}

现在, 我们证明 Binet--Cauchy 公式.

\begin{proof}
    设 \(B\) 的列~\(1\), \(2\), \(\dots\), \(m\)
    分别是
    \(b_1\), \(b_2\), \(\dots\), \(b_m\).
    设 \(A\) 的列~\(1\), \(2\), \(\dots\), \(n\)
    分别是
    \(a_1\), \(a_2\), \(\dots\), \(a_n\).
    那么
    \begin{align*}
        BA
        = {} & [Ba_1, Ba_2, \dots, Ba_n] \\
        = {} &
        \left[
        \sum_{k_1 = 1}^{m} {[A]_{k_1,1} b_{k_1}},
        \sum_{k_2 = 1}^{m} {[A]_{k_2,2} b_{k_2}},
        \dots,
        \sum_{k_n = 1}^{m} {[A]_{k_n,n} b_{k_n}}
        \right].
    \end{align*}
    利用行列式的多线性, 有
    \begin{align*}
             & \det {(BA)}
        \\
        = {} &
        \det {\left[
        \sum_{k_1 = 1}^{m} {[A]_{k_1,1} b_{k_1}},
        \sum_{k_2 = 1}^{m} {[A]_{k_2,2} b_{k_2}},
        \dots,
        \sum_{k_n = 1}^{m} {[A]_{k_n,n} b_{k_n}}
        \right]}
        \\
        = {} &
        \sum_{k_1 = 1}^{m} {[A]_{k_1,1}
        \det {\left[
        b_{k_1},
        \sum_{k_2 = 1}^{m} {[A]_{k_2,2} b_{k_2}},
        \dots,
        \sum_{k_n = 1}^{m} {[A]_{k_n,n} b_{k_n}}
        \right]}}
        \\
        = {} &
        \sum_{k_1 = 1}^{m} {
        \sum_{k_2 = 1}^{m} {
        [A]_{k_1,1} [A]_{k_2,2}
        \det {\left[
        b_{k_1},
        b_{k_2},
        \dots,
        \sum_{k_n = 1}^{m} {[A]_{k_n,n} b_{k_n}}
        \right]}}}
        \\
        = {} &
        \dots
        \\
        = {} &
        \sum_{1 \leq k_1, k_2, \dots, k_n \leq m}
        {[A]_{k_1,1} [A]_{k_2,2} \dots [A]_{k_n,n}
            \det {[b_{k_1}, b_{k_2}, \dots, b_{k_n}]}}.
    \end{align*}

    记
    \(
    f(k_1, k_2, \dots, k_n)
    = [A]_{k_1,1} [A]_{k_2,2} \dots [A]_{k_n,n}
    \det {[b_{k_1}, b_{k_2}, \dots, b_{k_n}]}.
    \)
    根据行列式的交错性,
    当 \(k_1\), \(k_2\), \(\dots\), \(k_n\) 中有二个数相同时,
    \(f(k_1, k_2, \dots, k_n) = 0\).

    (1)
    若 \(n > m\), 则 \(k_1\), \(k_2\), \(\dots\), \(k_n\)
    中总有二个相同
    (抽屉原理).
    所以, 每一项都是 \(0\).
    故 \(\det {(BA)} = 0\).

    (2)
    若 \(n \leq m\),
    那么 \(\det {(BA)}\) 等于
    \begin{align*}
        \sum_{1 \leq j_1 < j_2 < \dots < j_n \leq m}
        {\sum_{\substack{
        k_1, k_2, \dots, k_n\,\text{是} \\
                    j_1, j_2, \dots, j_n\,\text{的排列}
                }}
            {[A]_{k_1,1} [A]_{k_2,2} \dots [A]_{k_n,n}
                \det {[b_{k_1}, b_{k_2}, \dots, b_{k_n}]}}}.
    \end{align*}
    因为
    \(1 \leq j_1 < j_2 < \dots < j_n \leq m\),
    且
    \(k_1\), \(k_2\), \(\dots\), \(k_n\)
    是
    \(j_1\), \(j_2\), \(\dots\), \(j_n\)
    的排列,
    故
    \begin{align*}
        \det {[b_{k_1}, b_{k_2}, \dots, b_{k_n}]}
        = s(k_1, k_2, \dots, k_n)
        \det {[b_{j_1}, b_{j_2}, \dots, b_{j_n}]},
    \end{align*}
    从而 \(\det {(BA)}\) 等于
    \begin{align*}
         &
        \sum_{1 \leq j_1 < j_2 < \dots < j_n \leq m}
        {\det {[b_{j_1}, b_{j_2}, \dots, b_{j_n}]}}
        \\
         & \qquad \qquad \qquad
        \cdot
        \sum_{\substack{
        k_1, k_2, \dots, k_n\,\text{是} \\
                j_1, j_2, \dots, j_n\,\text{的排列}
            }}
        {s(k_1, k_2, \dots, k_n)\,
            [A]_{k_1,1} [A]_{k_2,2} \dots [A]_{k_n,n}}.
    \end{align*}
    注意,
    \begin{align*}
             &
        \sum_{\substack{
        k_1, k_2, \dots, k_n\,\text{是} \\
                j_1, j_2, \dots, j_n\,\text{的排列}
            }}
        {s(k_1, k_2, \dots, k_n)\,
            [A]_{k_1,1} [A]_{k_2,2} \dots [A]_{k_n,n}}
        \\
        = {} &
        \det {\left(
            A\binom{j_1, j_2, \dots, j_n}{1, 2, \dots, n}
            \right)}.
    \end{align*}
    再注意,
    因为
    \(1 \leq j_1 < j_2 < \dots < j_n \leq m\),
    故,
    由子阵的记号的定义,
    \begin{align*}
        [b_{j_1}, b_{j_2}, \dots, b_{j_n}]
        = B\binom{1, 2, \dots, n}{j_1, j_2, \dots, j_n}.
    \end{align*}
    从而
    \begin{align*}
             & \det {(BA)}
        \\
        = {} & \sum_{1 \leq j_1 < j_2 < \dots < j_n \leq m}
        {
            \det {\left(
                B\binom{1, 2, \dots, n}{j_1, j_2, \dots, j_n}
                \right)}
            \det {\left(
                A\binom{j_1, j_2, \dots, j_n}{1, 2, \dots, n}
                \right)}
        }.
    \end{align*}
\end{proof}

\end{document}

This material is about the (general) Binet--Cauchy formula
for the determinant of the product of two matrices.
It is of course optional.
The result is nontrivial,
and its several proofs (that I know) are not so easy.
The proof that I give here
uses the Leibniz formula (and some formulae related to it).
Another proof is recorded in the section
entitled "the Binet--Cauchy formula"
in appendix C,
which uses the generalised Laplace expansion
and two important properties of determinants:
\begin{itemize}
    \item replacing one column by
          the sum of that column
          and a multiple of another column
          keeps the determinant unchanged;
    \item the determinant of a square matrix with a zero row
          is zero.
\end{itemize}
